\chapter{Mastectomy}
\section*{Overview}
Mastectomy removes part or all of the breast, with the extent determined by the cancer, risk, and treatment plan.
\section*{Common Instruments and Equipment}
The exact setup varies by technique and setting. Common items may include:
\begin{itemize}[nosep,leftmargin=*]
  \item breast surgery instrument set
  \item electrosurgical or vessel-sealing device
  \item surgical drains and reconstruction supplies when applicable
\end{itemize}

\section*{Why It Is Done}
It may treat breast cancer or reduce cancer risk in people with a high inherited or familial risk.
\section*{How It Is Performed}
Options include total, skin-sparing, nipple-sparing, and modified radical mastectomy. Sentinel-node or axillary-node surgery and immediate reconstruction may be performed at the same operation.
\section*{Contraindications}
No single contraindication applies to every patient; the treating team assesses whether the procedure is suitable for the individual and the reason for treatment.

\section*{Precautions}
The care team reviews relevant health conditions, medicines, allergies, and procedure-specific risks, and explains any changes needed for the individual.

\section*{Risks and Recovery}
Bleeding, infection, seroma, numbness, arm swelling, shoulder stiffness, and body-image effects are possible. Drain care, pathology review, oncology planning, and reconstruction follow-up may be needed.
\section*{References}
\begin{enumerate}\item National Cancer Institute. Mastectomy. 2025.\item Current Medical Diagnosis and Treatment. McGraw-Hill; 2025.\end{enumerate}
\clearpage
